\section{Perfiles, evaluador y registro durable de cambios}\label{sec:matriz-perfiles-plazos}

La \cref{tab:matriz-perfiles-plazos} relaciona las fronteras descritas en las \cref{sec:implementacion-perfiles-plazos,sec:pruebas-perfiles-plazos}. La existencia de una prueba identifica su objeto de comprobación; los resultados de ejecución se conservan en \rutarepo{docs/verification-report.md}. No se suman los auxiliares ni las matrices de datos como campañas independientes.

\begin{table}[H]
  \centering
  \caption{Trazabilidad de perfiles y evaluación de plazos.}
  \label{tab:matriz-perfiles-plazos}
  \begin{tabularx}{\linewidth}{@{}L{3.2cm} X@{}}
    \toprule
    \textbf{Frontera} & \textbf{Comprobación discriminante} \\
    \midrule
    Solicitud de insumos & Reversibilidad de \texttt{DINP1}, referencia exacta, precisión original, ausencia frente a desconocimiento y rechazo de formas no canónicas. \\
    Definición & Reproducción del corpus, cantidades y corte independientes, referencias propias y conservación de calendarios completos. \\
    Canon de perfil & Lectura acotada de \texttt{DPRF1}, orden de colecciones, instante esperado con nanosegundos y desfase, y rechazo de truncamiento o exceso. \\
    Evaluador & Integridad antes de bloqueos, condiciones explícitas, cálculo parcial e instante final ausente cuando falta un dato necesario. \\
    Servicio y catálogo & Permisos, reautenticación, recibo \texttt{DPTX1}, revisión esperada, retiro terminal y aislamiento de colecciones. \\
    PostgreSQL & Atomicidad con auditoría, revalidación bajo bloqueo, carrera entre cambios y conservación de historia tras restaurar. \\
    Esquema e inventario & Tablas permanentes sin herencia, funciones y restricciones exactas, rol limitado, UUID cero y definición completa inválida detrás de una cabecera válida. \\
    Eventos & Fuente, revisión, operación y ámbito exactos; emisión transaccional, secuencia con huecos admitidos y restauración. \\
    Despacho & Selección por cabeza actual, familias y ámbito; trabajos, cursor y auditoría atómicos, con conciliación recurrente del legado. \\
    HTTP & Ruta y comando congruentes, entrada estricta, presupuesto del cuerpo y proyección completa de precisión y referencias. \\
    \bottomrule
  \end{tabularx}
\end{table}

{\raggedright
Las pruebas de aplicación están en el directorio
\rutarepo{crates/application/tests/}. Los archivos siguientes separan sus
fronteras:
\begin{itemize}
  \item \rutarepo{deadline_input_encoding.rs}: solicitud de insumos.
  \item \rutarepo{deadline_profile_encoding.rs}: canon del perfil.
  \item \rutarepo{deadline_profile_definition.rs}: definición y corpus.
  \item \rutarepo{deadline_profiled_evaluation.rs}: evaluación con perfil.
\end{itemize}
Los archivos cuyo nombre comienza con \texttt{deadline\_profile\_catalog},
en ese mismo directorio, cubren servicio, consultas, fallos y recibos.
\par}

{\raggedright
Las pruebas de infraestructura están en el directorio
\rutarepo{crates/infrastructure/tests/}. Para el catálogo comparten el prefijo
\texttt{deadline\_profile\_} y la extensión \texttt{.rs}. Los sufijos
\texttt{backend}, \texttt{revalidation}, \texttt{schema}, \texttt{restore},
\texttt{import} y \texttt{events} identifican sus fronteras.
El registro compartido se ejercita además en los siguientes archivos del mismo
directorio:
\begin{itemize}
  \item \rutarepo{deadline_source_events.rs}.
  \item \rutarepo{deadline_source_event_catalog.rs}.
  \item \rutarepo{deadline_source_event_restore.rs}.
\end{itemize}
Estas pruebas requieren servicios aislados configurados; su omisión por falta
de entorno no demuestra el adaptador.
\par}

La frontera de transporte usa los archivos con prefijo \texttt{deadline\_profile\_http} en \rutarepo{crates/web/tests/}. El contrato está en \rutarepo{docs/deadline-profiles-api.md}; la derivación del presupuesto JSON, en \rutarepo{docs/deadline-profile-json-budget.md}; y las obligaciones de respaldo y restauración, en \rutarepo{docs/database-operations.md}.

El seguimiento persistente y el consumidor local cuentan con las comprobaciones del \cref{sec:matriz-plazos-persistentes}. El servicio humano, HTTP y Qadra incorporan V2, y \texttt{serve} compone despacho y consumo con supervisión. Sus pruebas focales se registran separadas de las del catálogo y de los recorridos con servicios reales: la campaña de Qadra aprobó 25 escenarios, incluidos dos nuevos de seguimiento en escritorio y móvil; un recorrido HTTP comprobó reevaluación, señales, reinicios y restauración, mientras la repetición final de la campaña de API permanece pendiente. Agenda y alertas conservan su alcance pendiente. Esta matriz mantiene la separación entre una candidata matemática, una declaración de aplicabilidad y una actuación procesal acreditada; no sustituye ejemplos jurídicamente respaldados por la sola concordancia de emisores y lectores del prototipo.

\subsection{Asignación durable de trabajos}

\begingroup
\raggedright
El puerto de despacho está en \rutarepo{crates/application/src/deadline_dispatch/}. Su adaptador y la comprobación de apertura se separan en los módulos \texttt{deadline\_dispatch\_postgres} y \texttt{deadline\_dispatch\_schema} de \rutarepo{crates/infrastructure/src/}; las migraciones pertenecen al grupo \texttt{0019}. El contrato se documenta en \rutarepo{docs/deadline-dispatch.md}.

En \rutarepo{crates/infrastructure/tests/}, los archivos con prefijo \texttt{deadline\_dispatch\_} y extensión \texttt{.rs} separan las siguientes fronteras: \texttt{backend}, recorrido y reinicio; \texttt{families} y \texttt{candidates}, selección de dependencias; \texttt{atomicity} y \texttt{guards}, confirmación conjunta y restricciones; \texttt{catalog} e \texttt{inventory}, apertura e historia; \texttt{restore}, restauración; y \texttt{timeouts} y \texttt{measurement}, presupuestos por sentencia y mediciones. La evidencia de asignación de trabajos permanece separada de su conclusión por el consumidor y de cualquier entrega de avisos.
\par
\endgroup
